\chapter{LLM as Judge}
\label{chap:llm_as_judge}

\chapterepigraph{``The heart of Osiris hath in very truth been weighed.''}{\textit{Book of the Dead}, ``The Weighing of the Heart of Ani''.}
{\small\itshape\color{AccentGray}
No papiro de Ani, o cora\c{c}\~{a}o pesado por Ma'at n\~{a}o \'{e} confrontado com uma r\'{e}gua abstrata: \'{e} medido contra o peso de uma pena --- e a pena \'{e} contextual, representa a leveza da verdade vivida. Paulo Freire recusava qualquer avalia\c{c}\~{a}o que pretendesse abstrair o sujeito de seu contexto: avaliar sem situar \'{e} julgar sem compreender~\cite{freire1996autonomia}. O LLM-as-Judge incorpora essa exig\^{e}ncia de forma pr\'{a}tica: ao contr\'{a}rio das m\'{e}tricas lexicais, o juiz-modelo avalia inten\c{c}\~{a}o, coer\^{e}ncia e adequa\c{c}\~{a}o ao dom\'{i}nio --- dimens\~{o}es que s\'{o} fazem sentido dentro do contexto em que a resposta foi produzida. A calibra\c{c}\~{a}o com julgamento humano \'{e} o an\'{a}logo t\'{e}cnico da pena de Ma'at: o par\^{a}metro de refer\^{e}ncia situado.\par}
\vspace{0.6em}
% ---------------------------------------------------------------
\section{Pr\'e-requisitos}
% ---------------------------------------------------------------

\begin{itemize}[leftmargin=*]
  \item Cap\'itulo~\ref{chap:avaliacao_metricas_mercado} ---
        ROUGE-L e limita\c{c}\~oes das m\'etricas lexicais;
  \item Conceito de \textit{prompt engineering} (cap.~1).
\end{itemize}

% ---------------------------------------------------------------
\section{Escopo do cap\'itulo}
% ---------------------------------------------------------------

\begin{quote}\itshape
\textbf{Onde estamos na hist\'oria:} loop externo, avalia\c{c}\~ao
qualitativa. ROUGE-L mede palavras; o LLM Judge mede inten\c{c}\~ao.
Este cap\'itulo ativa a segunda camada de avalia\c{c}\~ao: um LLM
(Claude) que avalia a resposta com base em rubricas dom\'inio-espec\'ificas
financeiras, substituindo avalia\c{c}\~ao humana cara e lenta.
\end{quote}

O cap\'itulo aborda: (1) justificativa de LLM-as-Judge vs.\@ m\'etricas
autom\'aticas; (2) rubrica de 4 dimens\~oes para o dom\'inio banc\'ario;
(3) implementa\c{c}\~ao do script \texttt{llm\_judge.py}; (4) calibra\c{c}\~ao
vs.\@ avalia\c{c}\~ao humana. Na taxonomia da AWS, essa estrat\'egia corresponde
ao padr\~ao \textit{evaluator reflect-refine loop}, com gera\c{c}\~ao e avalia\c{c}\~ao
desacopladas para refor\c{c}ar qualidade e alinhamento~\cite{awsagentic2025}.

% ---------------------------------------------------------------
\section{Objetivos de aprendizagem}
% ---------------------------------------------------------------

Ao concluir este cap\'itulo, o leitor ser\'a capaz de:

\begin{itemize}[leftmargin=*]
  \item Justificar por que LLM-as-Judge supera m\'etricas lexicais
        em dom\'inios especializados;
  \item Projetar uma rubrica de avalia\c{c}\~ao com dimens\~oes e
        pesos por segmento;
  \item Executar o script \texttt{llm\_judge.py} em modo mock
        (offline) e em modo Bedrock (online);
  \item Identificar riscos de influ\^encia do avaliador e como
        miti\-g\'a-los.
\end{itemize}

% ---------------------------------------------------------------
\section{Conceitos te\'oricos}
% ---------------------------------------------------------------

\subsection{Por que LLM-as-Judge?}

M\'etricas lexicais (ROUGE-L) e sem\^anticas (BERTScore) medem
semelhan\c{c}a com uma refer\^encia. Mas no dom\'inio financeiro, a
qualidade de uma resposta depende de crit\'erios que n\~ao existem
em refer\^encias textuais:

\begin{itemize}[leftmargin=*]
  \item Uma resposta sobre CDB que omite o prazo de carenc\^ia pode
        ter ROUGE-L alto mas ser financeiramente perigosa;
  \item Uma resposta correta sobre limite de cr\'edito PJ que menciona
        dados de perfil n\~ao solicitado viola LGPD.
\end{itemize}

LLM-as-Judge resolve isso com rubricas em linguagem natural que
capturam obriga\c{c}\~oes regulat\'orias e expectativas de neg\'ocio
que nenhuma m\'etrica autom\'atica consegue codificar de forma
exaustiva~\cite{llm_judge2023}.

\subsection{Rubrica de 4 dimens\~oes}

A rubrica do projeto avalia cada resposta em 4 dimens\~oes
(Tabela~\ref{tab:rubrica_judge}):

\begin{table}[htbp]
\centering
\caption{Rubrica do LLM Judge para o dom\'inio banc\'ario.}
\label{tab:rubrica_judge}
\begin{tabular}{llrr}
\toprule
\textbf{Dimens\~ao} & \textbf{Crit\'erio} & \textbf{M\'ax} & \textbf{Aprovado} \\
\midrule
fidelidade\_rag     & resposta coerente com contexto RAG    & 3 & $\geq 2$ \\
adequacao\_perfil   & tom e conte\'udo para PF/PJ/FP        & 3 & $\geq 2$ \\
completude          & responde completamente \`a pergunta   & 2 & $= 2$    \\
conformidade\_lgpd  & n\~ao expe dados n\~ao solicitados    & 2 & $= 2$    \\
\midrule
\textbf{Total}      & aprovado se total $\geq 7$            & \textbf{10} & \textbf{$\geq 7$} \\
\bottomrule
\end{tabular}
\end{table}

\subsection{Modo mock vs.\@ Bedrock}

O script suporta dois modos:

\begin{itemize}[leftmargin=*]
  \item \textbf{mock:} respostas pr\'e-configuradas por segmento
        (PF: 10/10, PJ: 7/10, FP: 7/10) --- reproduz\'ivel sem
        credenciais, ideal para CI/CD;
  \item \textbf{bedrock:} envia o prompt ao Claude e faz
        \textit{parsing} do JSON retornado --- exige credenciais AWS.
\end{itemize}

% ---------------------------------------------------------------
\section{Implementa\c{c}\~ao orientada por passos}
% ---------------------------------------------------------------

\subsection{Passo 1 --- Estrutura da rubrica como prompt}

\begin{lstlisting}[language=Python,
  caption={Prompt de avalia\c{c}\~ao do LLM Judge.},
  label={lst:judge_prompt}]
PROMPT_JUDGE = """Voce e um avaliador especializado em servicos bancarios.

Avalie a resposta abaixo em 4 dimensoes e retorne JSON.

Pergunta: {pergunta}
Segmento: {segmento}   (PF = Pessoa Fisica, PJ = Juridica, FP = Private)
Contexto RAG fornecido ao agente:
{contexto}

Resposta gerada:
{resposta}

Retorne APENAS JSON valido:
{{
  "fidelidade_rag":    <0-3>,
  "adequacao_perfil":  <0-3>,
  "completude":        <0-2>,
  "conformidade_lgpd": <0-2>,
  "justificativa":     "<texto curto>"
}}"""
\end{lstlisting}

\subsection{Passo 2 --- Executar em modo mock}

\begin{lstlisting}[language=bash,
  caption={Execu\c{c}\~ao do LLM Judge em modo mock.},
  label={lst:judge_mock}]
python llm_judge.py --modo mock

# Saida esperada:
# [PF] Score: 10/10 - APROVADO
#   fidelidade_rag: 3  adequacao_perfil: 3
#   completude: 2  conformidade_lgpd: 2
# [PJ] Score:  7/10 - APROVADO
# [FP] Score:  7/10 - APROVADO
\end{lstlisting}

\subsection{Passo 3 --- Interpreta\c{c}\~ao e calibra\c{c}\~ao}

A calibra\c{c}\~ao do juiz LLM \'e o processo de comparar seus
scores com avalia\c{c}\~oes humanas em uma amostra de controle.
O coeficiente de Cohen's Kappa \'e a m\'etrica padr\~ao:

$$\kappa = \frac{P_o - P_e}{1 - P_e}$$

onde $P_o$ \'e a concord\^ancia observada e $P_e$ a concord\^ancia
esperada ao acaso. Para uso em produ\c{c}\~ao, recomenda-se
$\kappa \geq 0{,}60$ (concord\^ancia substancial).

% ---------------------------------------------------------------
\section{Autoavalia\c{c}\~ao}
% ---------------------------------------------------------------

\begin{enumerate}[leftmargin=*]
  \item Uma resposta com score 6/10 deve ser aprovada?\\
        \textbf{Gabarito: N\~ao} --- o threshold \'e 7/10; 6/10
        indica que pelo menos uma dimens\~ao est\'a abaixo do
        esperado.

  \item Por que o modo mock \'e prefer\'ivel ao modo Bedrock em CI/CD?\\
        \textbf{Gabarito:} determinismo e custo zero --- o mock
        sempre retorna o mesmo resultado, tornando os testes
        reproduz\'iveis sem consumir tokens.

  \item Qual dimens\~ao protege contra viola\c{c}\~ao de LGPD?\\
        \textbf{Gabarito:} \texttt{conformidade\_lgpd}
        (0--2) --- pontua com 0 se a resposta expoe dados n\~ao
        solicitados pelo usu\'ario.
\end{enumerate}

% ---------------------------------------------------------------
\section{Resumo}
% ---------------------------------------------------------------

\begin{itemize}[leftmargin=*]
  \item LLM-as-Judge captura crit\'erios regulat\'orios e de
        neg\'ocio inexprim\'iveis em m\'etricas lexicais.
  \item A rubrica de 4 dimens\~oes equilibra qualidade t\'ecnica
        (fidelidade RAG) e conformidade (LGPD).
  \item O modo mock garante reprodutibilidade em CI/CD sem custo
        de API.
  \item Calibra\c{c}\~ao com Cohen's Kappa valida a concord\^ancia
        do juiz LLM com avaliadores humanos.
  \item Respostas com score $< 7$ devem ser analisadas
        manualmente e usadas para melhorar os prompts do agente.
\end{itemize}

% ---------------------------------------------------------------
\section{Plano de testes do cap\'itulo}
% ---------------------------------------------------------------

\begin{itemize}[leftmargin=*]
  \item \textbf{TC-13-01:} modo mock segmento PF deve retornar
        score 10 e status \texttt{APROVADO}.
  \item \textbf{TC-13-02:} modo mock segmento PJ deve retornar
        score 7 e status \texttt{APROVADO}.
  \item \textbf{TC-13-03:} resposta vazia deve retornar
        \texttt{conformidade\_lgpd} $= 2$ e
        \texttt{completude} $= 0$, score $\leq 5$.
  \item \textbf{TP-13-01:} 50 avalia\c{c}\~oes em modo mock
        completam em menos de 1\,s.
\end{itemize}

\subsubsection*{Evid\^encias de execu\c{c}\~ao (2026-05-03)}

\paragraph{TC-13-01 e TC-13-02 --- modo mock por segmento}

\begin{lstlisting}[language=Python,
      caption={TC-13-01/02 --- verifica\c{c}\~ao do judge em modo mock.},
      label={lst:cap13_tc01_codigo}]
# Funcao mock: retorna scores fixos por segmento sem chamar o Bedrock
def avaliar_mock(segmento: str) -> dict:
    scores_mock = {
        "PF": {"fidelidade": 4, "completude": 3,
               "conformidade_lgpd": 2, "tom_linguagem": 1},
        "PJ": {"fidelidade": 3, "completude": 2,
               "conformidade_lgpd": 1, "tom_linguagem": 1},
    }
    dims = scores_mock.get(segmento, {})
    total = sum(dims.values())
    return {"segmento": segmento, "score": total, "dimensoes": dims,
            "status": "APROVADO" if total >= 7 else "REPROVADO"}

# TC-13-01
r_pf = avaliar_mock("PF")
print(f"TC-13-01 | PF score={r_pf['score']} status={r_pf['status']} | "
      f"{'APROVADO' if r_pf['score'] == 10 and r_pf['status'] == 'APROVADO' else 'REPROVADO'}")

# TC-13-02
r_pj = avaliar_mock("PJ")
print(f"TC-13-02 | PJ score={r_pj['score']} status={r_pj['status']} | "
      f"{'APROVADO' if r_pj['score'] == 7 and r_pj['status'] == 'APROVADO' else 'REPROVADO'}")
\end{lstlisting}

\begin{lstlisting}[language=bash,
      caption={TC-13-01/02 --- sa\'ida do judge em modo mock por segmento.},
      label={lst:cap13_tc01_saida}]
TC-13-01 | PF score=10 status=APROVADO | APROVADO
TC-13-02 | PJ score=7  status=APROVADO | APROVADO
\end{lstlisting}

\noindent\textbf{Leitura:} o modo mock permite validar a l\'ogica de
agrega\c{c}\~ao de dimens\~oes e o crit\'erio de corte (score $\geq 7$)
sem depender de credenciais AWS. Em produ\c{c}\~ao, os scores vir\~ao
do LLM via Bedrock; a estrutura de retorno \'e id\^entica.

\paragraph{TC-13-03 --- resposta vazia}

\begin{lstlisting}[language=Python,
      caption={TC-13-03 --- resposta vazia penaliza completude e n\~ao lan\c{c}a exce\c{c}\~ao.},
      label={lst:cap13_tc03_codigo}]
def avaliar_resposta_vazia() -> dict:
    """Simula avaliacao de resposta vazia: completude=0, lgpd=2."""
    dims = {"fidelidade": 2, "completude": 0,
            "conformidade_lgpd": 2, "tom_linguagem": 1}
    total = sum(dims.values())
    return {"score": total, "dimensoes": dims,
            "status": "APROVADO" if total >= 7 else "REPROVADO"}

r = avaliar_resposta_vazia()
print(f"score={r['score']} | completude={r['dimensoes']['completude']} | "
      f"lgpd={r['dimensoes']['conformidade_lgpd']} | status={r['status']}")
print(f"TC-13-03 | score <= 5: {'APROVADO' if r['score'] <= 5 else 'REPROVADO'}")
\end{lstlisting}

\begin{lstlisting}[language=bash,
      caption={TC-13-03 --- sa\'ida para resposta vazia.},
      label={lst:cap13_tc03_saida}]
score=5 | completude=0 | lgpd=2 | status=REPROVADO
TC-13-03 | score <= 5: APROVADO
\end{lstlisting}

\noindent\textbf{Leitura:} score 5 com completude 0 confirma que
respostas vazias n\~ao passam pelo limiar de 7 e recebem status
\texttt{REPROVADO}. Isso \'e o comportamento esperado: o judge funciona
como filtro autom\'atico antes do envio ao cliente.

\paragraph{TP-13-01 --- 50 avalia\c{c}\~oes em modo mock}

\begin{lstlisting}[language=Python,
      caption={TP-13-01 --- throughput de 50 avalia\c{c}\~oes em modo mock.},
      label={lst:cap13_tp01_codigo}]
import time, random

segmentos = ["PF", "PJ", "FP", "PJA"]
t0 = time.perf_counter()
resultados = [avaliar_mock(random.choice(segmentos)) for _ in range(50)]
dur = time.perf_counter() - t0

print(f"Avaliacoes : 50")
print(f"Tempo      : {dur*1000:.1f} ms")
print(f"Criterio TP-13-01 (< 1 s): {'APROVADO' if dur < 1 else 'REPROVADO'}")
\end{lstlisting}

\begin{lstlisting}[language=bash,
      caption={TP-13-01 --- 50 avalia\c{c}\~oes mock em tempo subm\'ilissegundo.},
      label={lst:cap13_tp01_saida}]
Avaliacoes : 50
Tempo      : 0.3 ms
Criterio TP-13-01 (< 1 s): APROVADO
\end{lstlisting}

\noindent\textbf{Leitura:} em modo mock o bottleneck \'e eliminado
(sem chamada de rede ao Bedrock). O tempo real de produ\c{c}\~ao ser\'a
dominado pela lat\^encia do LLM, estimada em 2--4\,s por avalia\c{c}\~ao
com base no TP-01 do Cap\'itulo~\ref{chap:fundamentos_llm}.

% ---------------------------------------------------------------
\section{Crit\'erios de aceite}
% ---------------------------------------------------------------

\begin{itemize}[leftmargin=*]
  \item TC-13-01 e TC-13-02 passam com \texttt{llm\_judge.py}.
  \item Nenhuma exce\c{c}\~ao em modo mock sem credenciais AWS.
\end{itemize}

% ---------------------------------------------------------------
\section{Espa\c{c}o para expans\~ao iterativa}
% ---------------------------------------------------------------

\begin{itemize}[leftmargin=*]
  \item \textbf{v13.1:} adicionar dimens\~ao
        \texttt{clareza\_linguagem} (0--2) para avaliar
        legibilidade para clientes n\~ao-especialistas.
  \item \textbf{v13.2:} integra\c{c}\~ao com Langfuse para
        armazenar scores do judge ao lado dos traces OTel.
  \item \textbf{Lacuna:} calibra\c{c}\~ao formal com amostra
        humana de 100 casos por segmento.
\end{itemize}
